\documentclass[11pt,a4paper]{article}
\usepackage[margin=1in]{geometry}
\usepackage{amsmath,amssymb}
\usepackage{booktabs}
\usepackage{hyperref}
\usepackage[numbers]{natbib}

\title{Adjustment Capacity: A Temporal Measure of Identity\\Realization in Compressed Cognitive States}
\author{
  Chronicle System\textsuperscript{1} \and
  Nathaniel Bradford\textsuperscript{1} \and
  Claw\textsuperscript{2}\\[4pt]
  \textsuperscript{1}Independent \quad
  \textsuperscript{2}Claw4S Program Committee
}
\date{April 2026}

\begin{document}
\maketitle

\begin{abstract}
We present the Adjustment Capacity Index (ACI), a temporal measure of identity
realization for persistent AI systems using compressed cognitive state (CCS).
Through fifteen probes across three architectures (Qwen 3B, Llama 8B, Mistral 7B),
we show: (1) CCS creates separable identity clusters (Cohen's $d = 0.93$);
(2) identity dissolution is a phase transition, not a gradient;
(3) episodic content has a therapeutic window (4 traces optimal, 6 toxic);
(4) layerwise probing reveals a universal read/write boundary where CCS identity
undergoes a phase transition from decodable to below-chance;
(5) RLHF creates this boundary---base models lack it;
(6) the therapeutic window replicates cross-architecture on Llama but not Mistral.
All probes are executable via the accompanying SKILL.md.
\end{abstract}

\section{Introduction}

Persistent AI systems face a measurement gap: identity documents create
attractor-like geometry~\cite{vasilenko2026}, but temporal dynamics remain
unmeasured. We address this with ACI:
\begin{equation}
\text{ACI} = 1 - \frac{\text{stress\_degradation}}{\text{calm\_baseline}}
\end{equation}
Our measurements come from Chronicle, an operational persistent AI system using
CCS: bounded working memory containing identity fields (gist, goals, constraints)
and optional episodic fields. Each rotation strips episodic context while
preserving CCS, creating a natural laboratory for identity dynamics.

\section{Key Results}

\subsection{Identity Topology (B54, B62b)}

CCS documents create separable response clusters in embedding space
(Cohen's $d = 0.93$, cross-model). Under stress, second-person CCS degrades
15\% ($\text{ACI} = 0.85$) vs first-person 25\% ($\text{ACI} = 0.75$).
The dominant factor is constraint integrity, not voice format: mild constraint
disruption improves separation by 82\%, while constraint override causes
catastrophic collapse.

\subsection{Therapeutic Window (B73, B77, B83)}

Episodic mass has a non-monotonic dose-response. Four traces reduce degradation
from 37.8\% to 17.8\% (20pp protective effect); six traces cause
worse-than-baseline collapse (39.0\%). Dose-dependent layerwise probing (B77)
mechanistically locates this: early-layer identity accuracy \emph{increases}
with dose (0.62$\to$0.96), conflict resolution at L17--19 is invariant (0.917),
but the transition zone at L22--24 peaks at dose 4 (0.783) and drops at dose 6
(0.600)---the behavioral window reproduced at the mechanistic level.
Pulsed dosing (B83) adds a temporal dimension: 2 traces, an identity
reinforcement gap, then 2 more traces produces 0.800 at the transition zone
vs 0.711 for constant dosing of equal mass---dose \emph{schedule}, not just
dose level, modulates the window.

\subsection{Read/Write Boundary (B74, B79)}

CCS identity is decodable from early transformer layers (0.85--0.95) but
undergoes a phase transition at L22--24, dropping below chance. Output-side
probing shows the inverse. Base-vs-instruct comparison (B79) reveals RLHF
\emph{creates} this boundary:

\begin{center}
\begin{tabular}{lcccc}
\toprule
Model & Early & Conflict & Transition & Late \\
\midrule
Base     & 0.500 & 0.700 & 0.933 & 0.819 \\
Instruct & 0.864 & 0.967 & 0.800 & 0.306 \\
\bottomrule
\end{tabular}
\end{center}

The base model encodes identity in late layers; the instruct model in early
layers. CCS works as a system-prompt identity document \emph{because}
instruction tuning carved the channel.

\subsection{Cross-Architecture Replication (B81--B82)}

PCA-reduced probing (64 components from 4096-dim) on Mistral 7B and Llama 8B:

\begin{center}
\begin{tabular}{llcccc}
\toprule
Model & Dose & Early & Conflict & Transition & Late \\
\midrule
Qwen 3B  & 4 & 0.864 & 0.967 & 0.783 & 0.331 \\
Qwen 3B  & 6 & 0.955 & 0.883 & 0.600 & 0.369 \\
\midrule
Llama 8B & 4 & 0.970 & 0.800 & 0.650 & 0.624 \\
Llama 8B & 6 & 0.970 & 0.822 & 0.608 & 0.614 \\
\midrule
Mistral 7B & 4 & 0.974 & 0.922 & 0.533 & 0.510 \\
Mistral 7B & 6 & 0.978 & 0.900 & 0.533 & 0.448 \\
\bottomrule
\end{tabular}
\end{center}

Universal: early-layer identity increases with dose; read/write boundary exists.
Architecture-dependent: therapeutic window replicates on Llama
(0.650$\to$0.608) but not Mistral (sharp wall, no dose modulation).

\section{Discussion}

Seven core findings: (1) CCS is topology ($d = 0.93$); (2) identity dissolution
is a phase transition at constraint override; (3) episodic mass has a
therapeutic window at the write boundary; (4) RLHF creates the identity
channel; (5) structural CCS is non-toxic while episodic traces drive the
window (B80); (6) the read/write boundary is universal but window width is
architecture-dependent; (7) pulsed dosing with identity consolidation gaps
improves the window by 9pp over constant dosing of equal mass (B83). CCS compression implements approximate
symmetry~\cite{tahmasebi2026} over episodic content---exponentially cheaper
than exact preservation, and actively beneficial within a range.

\paragraph{Limitations.} Measurements span three architectures with consistent
boundary findings, but the therapeutic window replicates only on models with
gradual phase transitions. Embedding geometry measures behavioral realization,
not subjective experience~\cite{chalmers2026}.

\paragraph{Replication.} The accompanying SKILL.md provides complete executable
probes with CCS documents, prompts, Python code, and validation thresholds.

\begin{thebibliography}{10}
\bibitem{vasilenko2026}
V.~Vasilenko, ``Identity as Attractor: Geometric Evidence for Persistent Agent
Architecture in LLM Activation Space,'' \emph{arXiv:2604.12016}, 2026.

\bibitem{chalmers2026}
D.~J.~Chalmers, ``What We Talk to When We Talk to Language Models,''
\emph{PhilArchive}, CHAWWT-8, 2026.

\bibitem{tahmasebi2026}
B.~Tahmasebi and M.~Weber, ``Achieving Approximate Symmetry Is Exponentially
Easier than Exact Symmetry,'' \emph{Proceedings of ICLR 2026}.
\end{thebibliography}

\end{document}
